\section{Models and the optimization problem of MPC}
\label{sec:problem}

\subsection{State-space models}

A \emph{dynamical system} is described by a set of \emph{state variables}, collected in the state
vector $x \in \R^n$, which summarize its past: given the state at one instant and the inputs from
then on, the future of the system is determined. The inputs are collected in $u \in \R^m$, the
measured outputs in $y \in \R^p$.

In continuous time the model is a differential equation, $\dot x(t) = f_c(x(t), u(t))$. The
controller, however, acts at the sampling instants, so the model it uses is in discrete time:
\begin{equation}
  x_{k+1} = f(x_k, u_k), \qquad y_k = h(x_k).
  \label{eq:model}
\end{equation}
The model is \emph{linear} when $f(x,u) = Ax + Bu$ and $h(x) = Cx$ for constant matrices $A$,
$B$, $C$; it is \emph{nonlinear} otherwise. Three examples were worked out in class: a mass
connected to a spring and a damper and pushed by a force, whose state is the pair (position,
velocity) and whose model is linear; a pendulum driven by a torque, nonlinear because of the
gravitational term $\sin\theta$; an RLC circuit, linear. The procedure is always the same: write
the physical balance equations, choose the states, the inputs and the outputs, and rewrite the
equations in the form \eqref{eq:model}.

An \emph{equilibrium} of \eqref{eq:model} under a constant input $\bar u$ is a state $x^{eq}$
that the system does not leave, $x^{eq} = f(x^{eq}, \bar u)$. A system may have one, several or
infinitely many equilibria, or none: the pendulum has two families of them, the hanging and the
upright position repeated at every full turn. Without loss of generality the equilibrium of
interest is placed at the origin, $f(0,0) = 0$, by a change of variables.

\subsection{The problem solved at every step}

The optimization problem that MPC solves at time $k$, from the measured state $x_k$, is
\begin{equation}
\begin{aligned}
  \mathcal{P}_N(x_k):\quad
  \min_{u_{\ik{0}},\dots,u_{\ik{N-1}}} \quad
    & V_N(x_k, \bu_k) = \sum_{i=0}^{N-1} \ell\big(x_{\ik{i}}, u_{\ik{i}}\big) + V_f\big(x_{\ik{N}}\big) \\
  \st \quad
    & x_{\ik{i+1}} = f(x_{\ik{i}}, u_{\ik{i}}), \qquad i = 0,\dots,N-1, \\
    & x_{\ik{0}} = x_k, \\
    & u_{\ik{i}} \in \U, \qquad i = 0,\dots,N-1, \\
    & x_{\ik{i}} \in \X, \qquad i = 0,\dots,N, \\
    & x_{\ik{N}} \in \Xf .
\end{aligned}
\label{eq:PN}
\end{equation}
Its ingredients are the following.
\begin{itemize}
  \item The \emph{horizon} $N$: the number of steps the controller looks ahead. A longer horizon
        gives a better plan at the price of a larger problem.
  \item The \emph{stage cost} $\ell(x,u)$: the cost paid at each step, which encodes the goal.
        For a regulation problem, where the state has to be brought to the origin, the usual
        choice is quadratic,
        \begin{equation}
          \ell(x,u) = x\T Q x + u\T R u,
          \label{eq:quadcost}
        \end{equation}
        with $R$ positive definite and $Q$ positive definite, or positive semi-definite when only
        some of the states are penalized.
  \item The \emph{terminal cost} $V_f(x)$: the cost charged on the last predicted state. Its role,
        which is central, is the subject of \cref{sec:optimality}.
  \item The \emph{dynamics} as equality constraints, with the measured state as initial condition:
        this is where the model enters, and what makes the controller a feedback.
  \item The sets $\U$ and $\X$ of admissible inputs and states: the limits of the process. The
        common case is a box, $u^{\min} \le u \le u^{\max}$, $x^{\min} \le x \le x^{\max}$.
  \item The \emph{terminal set} $\Xf$, a constraint on the last predicted state. It is not
        required by the application: it is added to guarantee that the problem remains solvable
        from one step to the next, as \cref{sec:feasibility} explains.
\end{itemize}
The optimal value of the problem is written $V_N^*(x_k)$, and the first element of the optimal
sequence defines the control law, $\kappa_N(x_k) = u^*_{\ik{0}}$.

\begin{example}[the tank]
A tank of water volume $x$ is filled by a pump of power $u$ and emptied by a known consumption
$w$. With a sampling time $T$, the volume evolves as $x_{k+1} = x_k + a\,u_k - b\,w_k$. The
volume has to stay between a minimum and a maximum, the pump between zero and its rated power,
and the level has to be kept close to a reference $x^{ref}$. The problem solved at every step is
\eqref{eq:PN} with $\ell(x,u) = (x - x^{ref})^2 + \rho\,u^2$, box constraints on $x$ and $u$, and
no terminal set. Because the model is linear and the cost quadratic, $\mathcal{P}_N(x_k)$ is a
quadratic program. The Tank Control Lab on the course website runs this controller next to a
PID on the same tank: the PID has no knowledge of the limits and of the future consumption,
and the difference is visible in the volume trajectory and in the number of minutes spent
outside the admissible band.
\end{example}

\subsection{The ideal problem: infinite horizon}

Behind \eqref{eq:PN} there is the problem one would really like to solve, the
\emph{infinite-horizon constrained optimal control problem}:
\begin{equation}
  \min_{u_0,u_1,\dots}\ \sum_{i=0}^{\infty} \ell(x_i, u_i)
  \quad \st\quad x_{i+1} = f(x_i,u_i),\ x_0 = x_k,\ u_i \in \U,\ x_i \in \X .
  \label{eq:infhor}
\end{equation}
It has infinitely many variables and constraints, and in general it cannot be solved. With a
linear model and the quadratic cost \eqref{eq:quadcost} it is the \emph{constrained LQR}
problem, still out of reach because of the constraints.

The one case that can be solved exactly is the \emph{unconstrained} linear-quadratic problem.
Its solution is a linear state feedback, $u^*_k = K x_k$, with the gain $K$ given by the
discrete-time Riccati equation, and the closed-loop cost from any initial state is a quadratic
form:
\begin{equation}
  J^*(x_0) = \sum_{k=0}^{\infty}\big(x_k\T Q x_k + u_k\T R u_k\big) = x_0\T W x_0 ,
  \label{eq:Wcost}
\end{equation}
where $W$ solves the Lyapunov equation of the closed loop. The derivation is short: under the
law $u = Kx$ the state is $x_k = (A+BK)^k x_0$, the input $u_k = K(A+BK)^k x_0$, and
substituting both in the sum collapses it into a single quadratic form in $x_0$. The matrix $W$
is the quantity that does all the work in the next chapter.

\begin{reading}
\refkey{Cannon\_2016}, section 2.1, ``Problem Description'' (pp.~13--15): the same problem and
the same notation as in these notes; section 2.2, ``The Unconstrained Optimum'' (pp.~15--18):
the cost under a linear feedback, the Lyapunov equation for $W$ and the Riccati equation for
the gain. \refkey{Rawlings\_2017}, section 2.2, ``Model Predictive Control'' (pp.~91--106), for
the general formulation with stage cost, terminal cost and terminal set, and section 1.2,
``Models and Modeling'' (pp.~1--10), for the modelling part.
\end{reading}
